% 판형: 152 x 225 mm (신국판)
\usepackage[paperwidth=152mm,paperheight=225mm,top=21mm,bottom=22mm,inner=19mm,outer=17mm,footskip=11mm,headsep=6mm]{geometry}
\usepackage{fontspec}
\usepackage{kotex}
\usepackage{amsmath}
\usepackage{xcolor}
\usepackage{graphicx}
\usepackage{array}
\usepackage{booktabs}
\usepackage{tabularx}
\usepackage{enumitem}
\usepackage{titlesec}
\usepackage{fancyhdr}
\usepackage[most]{tcolorbox}
\usepackage{tikz}
\usetikzlibrary{arrows.meta,positioning,shapes.geometric,shapes.misc,calc,decorations.pathreplacing,fit,backgrounds}
\usepackage{pgfplots}
\pgfplotsset{compat=1.18}
\usepackage[hidelinks]{hyperref}
\usepackage{xurl}

% 글꼴: 나눔 글꼴이 있으면 쓰고, 없으면 Windows의 맑은 고딕을 씁니다.
\IfFontExistsTF{NanumBarunGothic}{%
  \setmainhangulfont{NanumBarunGothic}[AutoFakeSlant=0.15]%
  \setmainfont{NanumBarunGothic}[AutoFakeSlant=0.15]%
  \setsansfont{NanumBarunGothic}%
}{%
  \setmainhangulfont{Malgun Gothic}[AutoFakeSlant=0.15]%
  \setmainfont{Malgun Gothic}[AutoFakeSlant=0.15]%
  \setsansfont{Malgun Gothic}%
}
\IfFontExistsTF{NanumSquareRound}{%
  \newfontfamily\titlefont{NanumSquareRound}[BoldFont={NanumSquareRound ExtraBold}]%
}{%
  \newfontfamily\titlefont{Malgun Gothic}%
}
\IfFontExistsTF{NanumGothicCoding}{%
  \setmonofont{NanumGothicCoding}%
}{%
  \setmonofont{Consolas}%
}

\linespread{1.5}
\setlength{\parindent}{1em}
\setlength{\parskip}{0.35em}
\frenchspacing
\emergencystretch=3em
% 한글에는 기울임꼴이 어울리지 않아서, 강조는 색으로 해요.
\renewcommand{\emph}[1]{\textcolor{inkblue}{#1}}
\newcommand{\filepath}[1]{{\small\nolinkurl{#1}}}

% 색
\definecolor{inkblue}{HTML}{2B5C8A}
\definecolor{skyfill}{HTML}{E8F1FA}
\definecolor{sunfill}{HTML}{FFF6D6}
\definecolor{sunline}{HTML}{E0B000}
\definecolor{leaffill}{HTML}{E9F6EA}
\definecolor{leafline}{HTML}{3C8D47}
\definecolor{rosefill}{HTML}{FCE9EC}
\definecolor{roseline}{HTML}{C0392B}
\definecolor{allowcol}{HTML}{2E86C1}
\definecolor{sendcol}{HTML}{E67E22}
\definecolor{greycol}{HTML}{8A8F98}

% 장 제목
\titleformat{\chapter}[display]
  {\titlefont\color{inkblue}}
  {\large\bfseries \thechapter 장}
  {0.6ex}
  {\bfseries\fontsize{22}{30}\selectfont}
  [\vspace{0.4ex}{\color{inkblue!40}\titlerule[1.2pt]}]
\titleformat{name=\chapter,numberless}[display]
  {\titlefont\color{inkblue}}
  {}
  {0pt}
  {\bfseries\fontsize{22}{30}\selectfont}
  [\vspace{0.4ex}{\color{inkblue!40}\titlerule[1.2pt]}]
\titlespacing*{\chapter}{0pt}{0pt}{16pt}
\titleformat{\section}{\titlefont\bfseries\large\color{inkblue}}{}{0pt}{}
\titlespacing*{\section}{0pt}{14pt}{6pt}
\titleformat{\subsection}{\titlefont\bfseries\normalsize\color{inkblue}}{}{0pt}{}
\titlespacing*{\subsection}{0pt}{10pt}{4pt}
\setcounter{secnumdepth}{0}
\setcounter{tocdepth}{0}

% 차례, 그림 이름
\renewcommand{\contentsname}{차례}
\renewcommand{\figurename}{그림}
\renewcommand{\tablename}{표}

% 쪽 머리와 쪽 번호
\pagestyle{fancy}
\fancyhf{}
\fancyhead[C]{\small\color{greycol}\leftmark}
\fancyfoot[C]{\small\thepage}
\renewcommand{\headrulewidth}{0pt}
\renewcommand{\chaptermark}[1]{\markboth{#1}{}}
\fancypagestyle{plain}{\fancyhf{}\fancyfoot[C]{\small\thepage}\renewcommand{\headrulewidth}{0pt}}

% 상자: 조금 더 알고 싶은 사람에게 (어른이나 고학년 독자용)
\newtcolorbox{deeper}[1][]{%
  enhanced, breakable, colback=skyfill, colframe=inkblue!70, boxrule=0.6pt, arc=2mm,
  left=3mm, right=3mm, top=2mm, bottom=2mm, fonttitle=\titlefont\bfseries\small,
  title={조금 더 알고 싶은 사람에게\ifx&#1&\else: #1\fi}, before skip=10pt, after skip=10pt,
  fontupper=\small}
% 상자: 생각해 보기
\newtcolorbox{think}{%
  enhanced, breakable, colback=sunfill, colframe=sunline, boxrule=0.6pt, arc=2mm,
  left=3mm, right=3mm, top=2mm, bottom=2mm, fonttitle=\titlefont\bfseries\small,
  coltitle=black, title={생각해 보기}, before skip=10pt, after skip=10pt}
% 상자: 이 장에서 알게 된 것
\newtcolorbox{learned}{%
  enhanced, breakable, colback=leaffill, colframe=leafline, boxrule=0.6pt, arc=2mm,
  left=3mm, right=3mm, top=2mm, bottom=2mm, fonttitle=\titlefont\bfseries\small,
  title={이 장에서 알게 된 것}, before skip=12pt, after skip=4pt}

\setlist[itemize]{leftmargin=1.4em, itemsep=0.15em, topsep=0.3em}
\setlist[enumerate]{leftmargin=1.6em, itemsep=0.15em, topsep=0.3em}

% 그림 공통 모양
\tikzset{
  kid/.style={circle, draw=inkblue, fill=skyfill, line width=0.7pt, minimum size=11mm, inner sep=0pt, font=\small},
  robot/.style={rounded rectangle, draw=greycol, fill=gray!12, line width=0.7pt, minimum height=8mm, inner sep=2pt, font=\small},
  fake/.style={circle, draw=roseline, dashed, fill=rosefill, line width=0.7pt, minimum size=9mm, inner sep=0pt, font=\scriptsize},
  allow/.style={-{Stealth[length=2.6mm]}, line width=1pt, color=allowcol},
  send/.style={-{Stealth[length=2.6mm]}, line width=1pt, color=sendcol},
  plainarrow/.style={-{Stealth[length=2.6mm]}, line width=0.9pt, color=inkblue},
}
\newcommand{\num}[1]{\textbf{#1}}
